\section{Introduction}

The adoption of automated decision-making systems in financial services has expanded rapidly, particularly in credit scoring and loan approval workflows. Machine learning models are increasingly used to assess creditworthiness based on historical applicant data. While these systems offer efficiency and scalability, they also inherit risks from the data on which they are trained. If sensitive or socially correlated features such as gender, age, or nationality are present—either directly or through proxies—the model may produce outcomes that disproportionately disadvantage certain groups.

This concern has brought algorithmic fairness to the forefront of AI governance and risk management. In credit decision contexts, unfair treatment can lead not only to reputational and ethical issues, but also to regulatory non-compliance. Traditional approaches to fairness in machine learning often focus on removing protected attributes from the dataset (fairness through unawareness), but this is insufficient when bias persists via correlated features. To address this, the field has developed a range of fairness-aware methods for auditing and mitigating bias throughout the model lifecycle. However, many practical implementations—especially in financial applications—lack accessible, reproducible workflows that illustrate these methods using real-world data.


